\section{Kesimpulan}

Algoritma \textit{Neural Style Network} (NST) yang nampak kompleks dan menggunakan istilah yang rumit ternyata hanya merupakan penerapan konsep-konsep Aljabar Linier tingkat tinggi. Dalam sintesis citra artistik, aljabar linier yang digunakan mulai dari Matriks Gram dan \textit{Convolutional Neural Network} (NST) sebenarnya hanya hasil abstraksi dari konsep-konsep aljabar linier dasar. Representasi citra digital dapat dipandang sebagai sebuah matriks dan/atau tensor berdimensi tinggi, telah dibahas bahwa setiap citra dapat dipandang sebagai elemen dalam ruang vektor sehingga operasi untuk aljabar vektor dapat diterapkan. Lalu, setelah dipetakan/direpresentasikan menggunakan vektor dan tensor, CNN dapat digunakan untuk memetakan ruang piksel menjadi ruang fitur, di mana setiap layer menghasilkan \textit{feature map} yang merepresentasikan citra pada layer abstraksi yang berbeda.

Dalam konteks NST, citra utama yang dipertimbangkan adalah citra konten $I_c$, citra gaya $I_s$, dan citra hasil $I_g$ yang digunakan sebagai variabel optimasi. Ketiganya dimasukkan ke dalam jaringan CNN VGG-19 sehingga pada leyer-layer tertentu diperoleh matriks fitur $F_c^l, F_s^l, F_g^l$ yang membaca konten dan gaya dalam basis yang sama. Di mana jaringan CNN ini sebenarnya hanyalah transformasi linier (yang dapat direpresentasikan oleh sebuah matriks) yang dilakukan dalam beberapa layer.

Representasi konten didefinisikan sebagai pola aktivasi pada \textit{feature map}, sehingga \textit{content loss} dapat dihitung dengan menghitung jarak kuadrat (norma Frobenius) antara $F_c^l$ dan $F_g^l$ untuk memastikan citra hasil mendekati citra konten. Sedangkan, representasi gaya diperoleh melalui Matriks Gram $G_s^l$ dan $G_g^l$ yang dibentuk dari \textit{inner product} antar fitur pada beberapa layer gaya. Matriks Gram dapat menangkan korelasi statistik antar fitur dan bersifat simetris, sehingga dapat dipandang sebagai matriks korelasi berisikan tekstur, pola, dan palet warna dari sebuah citra gaya, tanpa bergantung pada posisi aktual dari detail gaya tersebut. Sama, seperti pada \textit{content loss}, \textit{style loss} dapat dihitung menggunakan jarak kuadrat antara Matriks Gram dari citra hasil dan Matriks Gram dari citra gaya.

Sementara itu, representasi gaya diperoleh melalui Matriks Gram $G_s^l$ dan $G_g^l$ yang dibentuk dari inner product antar kanal fitur pada beberapa layer gaya. Matriks Gram ini menangkap korelasi statistik antar fitur dan bersifat simetris serta positive semi-definite, sehingga dapat dipandang sebagai matriks korelasi yang menyandikan tekstur, pola, dan palet warna citra gaya tanpa bergantung pada posisi spasial detail tersebut. Style loss kemudian dirumuskan sebagai jarak antara Gram matrix citra hasil dan citra gaya, sehingga meminimalkan style loss akan memaksa korelasi fitur $I_g$ mengikuti pola korelasi $I_s$.

Terakhir, kedua komponen loss tersebut digabungkan dalam persamaan \textit{loss total}, 
$$
L_{\text{total}}(I_g) = \alpha L_{\text{content}}(I_g) + \beta L_{\text{style}}(I_g)
$$

dengan hiperparameter $\alpha$ dan $\beta$ yang mengatur trade-off intensitas gaya dan konten. Proses optimasi dilakukan secara iteratif dengan menghitung gradien $\nabla_{I_g} L_{\text{total}}$ terhadap piksel $I_g$ melalui aturan rantai dan backpropagation, kemudian memperbarui $I_g$ sedikit demi sedikit dengan skema optimasi berbasis gradien hingga diperoleh citra hasil yang secara visual menggabungkan konten $I_c$ dan gaya $I_s$.

Dari sudut pandang Aljabar Linier, keseluruhan algoritma NST menunjukkan bagaimana representasi citra sebagai tensor, operasi matriks (termasuk transpos, perkalian, dan pembentukan Matriks Gram), konsep inner product dan norma, serta perhitungan gradien berperan langsung dalam memformulasikan dan menyelesaikan masalah pemisahan dan penggabungan konten-gaya pada citra digital. Dengan demikian, NST dapat dipandang sebagai contoh nyata bagaimana teori ruang vektor dan transformasi linier diaplikasikan untuk menghasilkan sintesis citra artistik yang kompleks namun tetap didasarkan pada struktur matematis yang jelas dan terukur.

